\subsection{两个集合的乘积}

定理1. 关系

\[
(\forall \mathrm{X}) (\forall \mathrm{Y}) (\exists \mathrm{Z}) (\forall z) ((z \in \mathrm{Z}) \iff (\exists x) (\exists y) (z = (x, y) \text {   且   } x \in \mathrm{X} \text {   且   } y \in \mathrm{Y})
\]

是真的。换句话说，对于所有 X 和所有 Y，关系“z 是一个有序对且 \(pr_{1}z \in X\) 并且 \(pr_{2}z \in Y\) ”在 z 上是集合化的。

让 \(A_y\) 表示所有形如 \((x, y)\) 为了 \(x \in X\) （参见§1，第6号，准则C53）。设 \(R\) 为关系 \(z \in A_y\) ，它等价于 \((\exists x)(z = (x, y)\) 并且 \(x \in X)\) 。显然，这个关系

\[
(\forall y) (\exists \mathrm{A}) (\forall z) (\mathrm{R} \Longrightarrow (z \in \mathrm{A}))
\]

是真的，依据S5（第一章，§4，第2节）成立。随后由S8可知，关系 \((\exists y)(y \in \mathbf{Y}\) 并且 \(\mathbf{R})\) 关于 \(z\) 是集合化的。但此关系等价于 \((\exists x)(\exists y)(y \in \mathbf{Y}\) 并且 \(x \in \mathbf{X}\) 并且 \(z = (x, y)\) ）；因此结论成立。

定义1. 给定两个集合X和Y，集合

\[
\mathcal {E} _ {z} ((\exists x) (\exists y) (z = (x, y) \text {   且   } x \in X \text {   且   } y \in Y))
\]

称为X与Y的积，并记为 \(X \times Y\) .

关系 \(z \in X \times Y\) 因此等价于“ \(z\) 是一个有序对，且 \(\text{pr}_1z \in X\) 并且 \(\text{pr}_2z \in Y\) ''。集合 \(X\) 和 \(Y\) 被称为 \(X \times Y\) 的第一因子和第二因子.

命题1. 若A'、B'为非空集合，则关系A' × B' ⊂ A × B等价于“A' ⊂ A且B' ⊂ B”。

首先，这种关系 \(z \in A' \times B'\) 等价于“ \(z\) 是一个有序对，且 \(\mathrm{pr}_1z \in A'\) 并且 \(\mathrm{pr}_2z \in B'\)''；因此，在没有任何限制的情况下， \(A'\) 并且 \(B'\) ，关系“ \(A' \subset A\) 并且 \(B' \subset B\)'' 蕴含

\[
\mathbf {A} ^ {\prime} \times \mathbf {B} ^ {\prime} \subset \mathbf {A} \times \mathbf {B}.
\]

反之，先证明当 \(B' \neq \emptyset\) 时（不对 \(A'\) 作限制），关系 \(A' \times B' \subset A \times B\) 蕴含 \(A' \subset A\) 。设 x 是 \(A'\) 的一个元素；由于 \(B' \neq \emptyset\) ，存在属于 \(B'\) 的对象 y；于是 \((x, y) \in A' \times B'\) ，从而 \((x, y) \in A \times B\) ，故 \(x \in A\) ；因此 \(A' \subset A\) . 类似地，当 \(A' \neq \emptyset\) 时，关系 \(A' \times B' \subset A \times B\) 蕴含 \(B' \subset B\) 。结论得证。

命题2. 设 A 和 B 是两个集合。关系 \(\mathbf{A} \times \mathbf{B} = \varnothing\) 等价于“ \(\mathbf{A} = \varnothing\) 要么 \(\mathbf{B} = \varnothing\) ''.

关系 \(z \in A \times B\) 蕴含 \(\mathsf{pr}_1z \in A\) 且 \(\mathsf{pr}_2z \in B\) ；因此 \(A \neq \emptyset\) 且 \(B \neq \emptyset\) 。反之，关系“ \(x \in A\) 且 \(y \in B\) ”蕴含 \((x, y) \in A \times B\) ，从而 \(A \times B \neq \emptyset\) 。换言之，关系 \(A \times B \neq \emptyset\) 等价于“ \(A \neq \emptyset\) 且 \(B \neq \emptyset\) ”；结论得证。

若 A、B、C 是集合，我们记

\[
(\mathbf {A} \times \mathbf {B}) \times \mathbf {C} = \mathbf {A} \times \mathbf {B} \times \mathbf {C};
\]

一个元素 \(((x, y), z)\) 的 A × B × C 写作 \((x, y, z)\) 并称为三元组。同样，若 A、B、C、D 是集合，我们写作

\[
(\mathbf {A} \times \mathbf {B} \times \mathbf {C}) \times \mathbf {D} = \mathbf {A} \times \mathbf {B} \times \mathbf {C} \times \mathbf {D};
\]

等等。

